\documentclass[11pt, letterpaper]{article}

\usepackage[top=1in, bottom=1in, left=1in, right=1in]{geometry}
\usepackage[T1]{fontenc}
\usepackage[utf8]{inputenc}
\usepackage{amsmath, amssymb, amsfonts}
\usepackage{newtxtext}
\usepackage{newtxmath}
\usepackage{microtype}
\usepackage{bm}
\usepackage{graphicx}
\usepackage[font=small, labelfont=bf]{caption}
\usepackage{subcaption}
\graphicspath{{./figures/}}
\usepackage{booktabs}
\usepackage{array}
\usepackage{multirow}
\usepackage{tabularx}
\usepackage[dvipsnames]{xcolor}
\usepackage[colorlinks=true, linkcolor=NavyBlue, citecolor=NavyBlue, urlcolor=NavyBlue]{hyperref}
\usepackage[numbers, sort&compress]{natbib}
\usepackage{enumitem}
\usepackage{setspace}
\usepackage{fancyhdr}
\usepackage{lastpage}
\usepackage{titlesec}
\usepackage{xspace}
\usepackage{float}
\usepackage{algorithm}
\usepackage{algpseudocode}
\usepackage{siunitx}

\setstretch{1.09}

\setlength{\abovedisplayskip}{8pt plus 2pt minus 4pt}
\setlength{\belowdisplayskip}{8pt plus 2pt minus 4pt}
\setlength{\abovedisplayshortskip}{4pt plus 2pt minus 2pt}
\setlength{\belowdisplayshortskip}{4pt plus 2pt minus 2pt}

\setlength{\parindent}{0.5in}
\setlength{\parskip}{0pt}

\titleformat{\section}{\normalfont\Large\bfseries}{\thesection.}{0.5em}{}
\titleformat{\subsection}{\normalfont\large\bfseries}{\thesubsection}{0.5em}{}
\titleformat{\subsubsection}{\normalfont\normalsize\bfseries}{\thesubsubsection}{0.5em}{}
\titlespacing*{\section}{0pt}{8pt plus 2pt minus 2pt}{3pt plus 1pt minus 1pt}
\titlespacing*{\subsection}{0pt}{5pt plus 2pt minus 1pt}{2pt plus 1pt minus 1pt}
\titlespacing*{\subsubsection}{0pt}{3pt plus 1pt minus 1pt}{1pt plus 1pt minus 1pt}

\pagestyle{fancy}
\fancyhf{}
\renewcommand{\headrulewidth}{0pt}
\fancyfoot[C]{\thepage}

\newcommand{\proposaltitle}{Advancements for Image Classification in Real-World Settings: Innovations in Distributed Dataset Acquisition and Cost-Sensitive Machine Learning}
\newcommand{\studentname}{Zi K. Deng}
\newcommand{\advisorname}{Jeffrey J. Rodriguez}
\newcommand{\advisortitle}{Associate Professor}
\newcommand{\departmentname}{Department of Electrical \& Computer Engineering}
\newcommand{\universityname}{The University of Arizona}
\newcommand{\proposaldate}{\today}
\newcommand{\ie}{i.e.,\xspace}
\newcommand{\eg}{e.g.,\xspace}
\newcommand{\etal}{et~al.\xspace}

\begin{document}

\begin{titlepage}
\centering
\vfill
{\LARGE \textbf{ECE Written Comprehensive Exam\\[6pt] Research Proposal}}
\vfill
{\Large \textbf{\proposaltitle}}
\vfill
{\large \studentname \\[6pt] \departmentname \\[4pt] \universityname \\[4pt] Tucson, Arizona}
\vfill
{\large \textbf{Faculty Advisor}: \advisorname \\[4pt] \advisortitle, \departmentname}
\vfill
{\large \textbf{Comprehensive Exam Committee:}\\[6pt]
Jerzy W. Rozenblit, Professor of Electrical \& Computer Engineering \\[4pt]
Abhijit Mahalanobis, Associate Professor of Electrical \& Computer Engineering \\[4pt]
Chicheng Zhang, Assistant Professor of Computer Science}
\vfill
{\large \proposaldate}
\vfill
\end{titlepage}

\section*{Abstract}
\addcontentsline{toc}{section}{Abstract}

Deep learning models for image classification have achieved remarkable
performance on curated benchmarks, yet deploying these models in applied
domains such as biodiversity monitoring and medical screening exposes two significant gaps. First, constructing training datasets from
distributed web sources often requires downloading images from multiple
heterogeneous servers. Many of these servers enforce rate limits that
existing tools neither detect nor respect, which can lead to systematic download
failures and biased training data. Second, standard classification objectives
treat all misclassification errors as equally costed. However, specific errors often carry consequences that differ by orders
of magnitude, such as misidentifying a dangerous species as harmless or
undergrading a severe disease stage. This proposal presents two contributions addressing
these gaps.

In our first contribution, we develop FLOW-DC (Flexible Large-scale
Orchestrated Workflow for Data Collection), a distributed image downloading
framework for constructing ML-ready datasets from URL manifests hosted across
heterogeneous servers. Existing tools such as img2dataset lack
adaptive rate-control mechanisms: they treat server-side throttling as
permanent failure, leading to download failures concentrated on the
small institutional servers that lack large-scale infrastructure. To address this, FLOW-DC introduces PAARC
(Policy-Aware Adaptive Request Controller), a BBR (Bottleneck Bandwidth and Round-trip propagation time) inspired application-layer
congestion control algorithm that dynamically adjusts per-host request
concurrency based on TTFB (time to first byte) latency, HTTP status codes, and connection-level
error signals. PAARC operates within the TaskVine distributed workflow
manager, enabling elastic execution across heterogeneous compute
environments (university HPC clusters, cloud instances, and local
workstations) without requiring a shared filesystem or synchronized start
times. FLOW-DC has been deployed at production scale for the construction of
the 38.7-million-image Biotrove dataset and will be evaluated against
img2dataset on both content delivery network-based and server-diverse workloads using
server-side failure rate, scaling efficiency, and class distribution fidelity
as the primary metrics.

In our second contribution, we address the inability of standard image classifiers to distinguish between misclassification errors of differing cost. In many applied domains, certain prediction errors are far more consequential than others. For example, misidentifying a deadly spider as harmless could delay life-saving treatment, while the reverse error results only in unnecessary caution. Standard cross-entropy training treats all errors as equally costly and provides no mechanism for a domain expert to encode their knowledge on misclassification costs. To address this, we present NICME (Non-Identical Cost Matrix Embedding), a cost-sensitive classification framework that accepts a user-defined misclassification cost matrix and integrates it into both the logit computation and a dedicated regularization term. NICME utilizes LoRA (Low-Rank Adaptation) for fine-tuning to prevent the model from memorizing the training data entirely. NICME will be evaluated across four datasets spanning binary classification to 20-class identification tasks in domains including spider species classification, diabetic retinopathy grading, and pharmacy medication dispensing. It will be compared against standard cross-entropy, as well as the regularization and logit adjustment components applied individually.

\vspace{12pt}
\noindent\textbf{Keywords:} cost-sensitive learning,
congestion control, image classification, deep learning,
cost matrix, distributed download

\newpage

\section{Research Topic and Its Importance}
\subsection{Introduction}
\label{sec:research_topic}
Deep learning models for computer vision have improved at a rapid pace over the past two decades, progressing from the convolutional architectures that first achieved production-grade image classification such as AlexNet and ResNet~\cite{krizhevsky2012,he2016resnet}, to the dominance of transformer-inspired architectures such as ViT, Swin, and ConvNeXt~\cite{dosovitskiy2021image, liu2021swin, liu2022convnext}, and now to the current ecosystem of compute-intensive self-supervised foundation models such as MAE, CLIP, and DINO~\cite{he2022mae, radford2021clip, darcet2025dinov3}. These advances have been matched by a corresponding growth in dataset scale: from CIFAR-10's 60{,}000 images totaling approximately 170~MB~\cite{krizhevsky2009} to LAION-5B's 5.85 billion image-text pairs totaling approximately 220~TB~\cite{schuhmann2022laion5b}. Each generation of models depends on increasingly large and more diverse training corpora to achieve better performance and stronger generalization.

However, the infrastructure and assumptions underlying these benchmark datasets do not reflect the conditions that researchers commonly encounter when building classification systems for applied domains. Benchmark datasets are distributed as self-contained archives with uniform formatting, generally balanced or well-characterized class distributions, and have their images available from high-capacity servers and fast networks~\cite{torralba2011unbiased}. In contrast, constructing a training dataset for domains such as biodiversity monitoring or practical medical screening often requires downloading millions of images from heterogeneous web sources, such as university servers, museum repositories, and citizen-science platforms~\cite{yang2024biotrove}. Many of these resources operate on limited server infrastructure and enforce rate limits that existing download tools neither detect nor respect~\cite{deng2025flowdc}. Once such a dataset is assembled, another gap often emerges at the model training stage: standard classification objectives implicitly assume that all misclassifications are equally damaging, an assumption that breaks down in applied settings where some errors carry severe real-world consequences. Misidentifying a dangerous species as harmless carries a significantly different implication than confusing two harmless species~\cite{deng2026aranet}; similarly, failing to detect a severe stage of diabetic retinopathy carries different consequences than overgrading a mild case~\cite{galdran2020}. The theoretical foundations for incorporating such asymmetric costs into classification have existed for decades~\cite{elkan2001}, yet their integration into modern deep learning architectures remains limited.

This research addresses both gaps through two complementary lines of investigation. First, we have developed a distributed image downloading pipeline that utilizes an application-layer congestion control algorithm to adapt request rates to individual host server capacity and allowances, enabling researchers and groups operating outside of resource-rich organizations to construct large-scale training datasets from server-diverse web sources in a timely manner without triggering rate-limit violations. Second, we are developing a cost-sensitive learning (CSL) framework that accepts a user-defined $N \times N$ misclassification cost matrix as a training hyperparameter, incorporated into a training pipeline that utilizes both cost-sensitive regularization and logit adjustment. We believe the combination of a robust data downloader and a CSL framework decoupled from class imbalance will provide meaningful improvements to practical image classification and strengthen the end-to-end machine learning (ML) pipeline across a variety of application domains.

\subsection{Limitations of Current Cost-Sensitive Learning and Dataset Acquisition Paradigms}
\label{sec:objectives}
Our research has identified two specific technical problems that arise when tackling image classification pipelines for real-world applications. The first is the lack of a dataset construction tool that can utilize adaptive per-host congestion control to mitigate server overload and throttling. To be an improved alternative to leading dataset downloaders, such a tool must be able to execute in a distributed manner, ideally across heterogeneous compute environments to accommodate varying infrastructure constraints. The problem with existing tools is that they lack adaptive rate-control mechanisms, treating server-side throttling as permanent failures, and require homogeneous compute cluster infrastructure such as Apache Spark~\cite{zaharia2016spark, beaumont2022img2dataset} that is either unavailable or impractical to set up in many academic settings. The second is the absence of cost-sensitive deep learning methods that operate with user-defined cost matrices on modern vision architectures. Existing approaches either derive cost weights from class frequency statistics rather than domain knowledge~\cite{menon2021logitadjust, khan2018costsensitive}, or have only been validated on outdated network architectures~\cite{chung2016csdnn}. Accordingly, our research objectives are the following:

\paragraph{\textbf{Objective 1: Distributed Downloading with Adaptive Congestion Control.}} Our first objective was to develop a distributed image download system that achieves download throughput comparable to the current state-of-the-art distributed downloaders, but significantly reduces the failure rate from server-side faults such as HTTP 429, 5XX, and connection reset errors.

\paragraph{\textbf{Objective 2: Cost-Sensitive Classification with User-Defined Cost Matrices.}} Our second objective is to develop a cost-sensitive image classification framework that aims to attain higher recall on user-designated high-cost classes than standard cross-entropy loss training and rivals existing cost-sensitive methods. To achieve this, we specifically introduce an $N \times N$ cost matrix as a hyperparameter to be utilized during model training.

\section{Review of Current Technology and Motivation}

\subsection{Distributed Image Dataset Construction}

The growth of web-scale vision datasets has created a need for tools that can
convert URL manifests, tabular files listing the remote locations of millions
of images, into locally stored, ML-ready training sets. Two systems have
emerged as the primary tools for this task: img2dataset, which underpins the
construction of most large-scale web-crawled vision datasets, and the Imageomics
distributed-downloader, which was developed specifically for biodiversity image
acquisition from institutional repositories.

The tool img2dataset, created by Beaumont and central to the construction of LAION-400M,
LAION-5B, and DataComp, is the most widely adopted open-source tool for
large-scale image downloading and
packaging~\cite{schuhmann2022laion5b,beaumont_img2dataset}. Its architecture
follows a staged pipeline: a reader ingests URL lists from one of thirteen
supported input formats via PyArrow, a distributor dispatches shards to a pool
of worker processes, and within each process a thread pool of up to 256 threads
performs concurrent HTTP downloads using Python's synchronous
\texttt{urllib.request} library~\cite{beaumont_img2dataset}. Downloaded images
are resized via OpenCV and written to one of five output formats, with
WebDataset tar archives recommended for datasets exceeding one million samples. With optimized network speeds, img2dataset achieves approximately 90 samples per second per CPU core, scaling
to roughly 1,300 images per second on a 16-core machine and approximately
10,000 images per second across a 10-node Spark
cluster~\cite{schuhmann2022laion5b}. The tool's largest documented deployment
downloaded LAION-5B's 5.85 billion image-text pairs using ten AWS c6i.4xlarge
instances over approximately one week, producing 220~TB of output at 384-pixel
resolution.

A notable architectural characteristic of img2dataset is its lack of any
congestion control or host-aware request scheduling. All URLs within a shard are
downloaded as fast as the thread pool allows, with no throttling based on server
response times, HTTP status codes, or per-domain rate limits. When a server responds with HTTP~429
(Too Many Requests), img2dataset treats this as a standard failure, no different than if the content doesn't exist; by
default, no retries are attempted (\texttt{retries=0}). This design is well
suited to downloading from content delivery networks (CDNs) such as Cloudflare and AWS
CloudFront, which tolerate high concurrency without rate limiting. However, when
the target data resides on smaller institutional servers, as is characteristic
of biodiversity repositories, museum collections, and regional medical
archives, the absence of rate control produces cascading failures: excessive
rate-limiting errors can lead to permanently failed downloads, and in severe cases,
IP-level blocks from the host server.

The Imageomics distributed-downloader, developed by the TreeOfLife project for
constructing biodiversity datasets from GBIF (Global Biodiversity Information
Facility) occurrence records, addresses a different deployment
scenario~\cite{yang2024biotrove}. This tool is designed around Slurm-based
queue submission with MPI (message passing interface) for inter-node coordination, making it well suited to
university HPC clusters where Slurm is the standard job scheduler. The
distributed-downloader was used to construct portions of the TreeOfLife dataset,
which aggregates species images from across GBIF's network of over 500
contributing institutional servers. Unlike img2dataset,
the distributed-downloader operates in an environment where the target servers
are heterogeneous in capacity, making the presence of per-host rate
control a necessity.

However, the distributed-downloader shares several limitations with img2dataset.
It does not implement per-host congestion control or adaptive rate pacing;
downloads proceed at the maximum rate the worker can sustain regardless of
server load signals. Its reliance on MPI requires a fixed world size determined
at job submission time; workers cannot join or leave the workflow dynamically,
and the compute resources must belong to a single Slurm-managed cluster. This
excludes scenarios in which a research group might combine university HPC
allocations, NSF cloud instances such as Jetstream2, and local workstations into
a single coordinated download workflow.

The limitations of both systems converge on two gaps. First, neither tool
implements adaptive per-host request rate control informed by server response
signals such as TTFB latency, HTTP~429 and 503 codes, or connection-level
errors. Second, neither supports distributed execution across heterogeneous
compute environments without requiring a shared filesystem, a common job
scheduler, or synchronized worker start times. The practical consequence of
these gaps is that academic research groups attempting to construct training
datasets from server-diverse sources, the typical case for biodiversity,
ecology, and regional medical imaging, face systematic download failures
concentrated on the smallest institutional servers. Because these servers
disproportionately host images from rare or geographically restricted taxa, the
resulting dataset underrepresents precisely the classes that are most
ecologically or clinically relevant, possibly introducing a bias into the training data
before any model is trained.

\subsection{Cost-Sensitive Deep Learning}

Although cost-sensitive classification and class-imbalanced classification are
frequently treated as interchangeable in the deep learning literature, they are
mathematically distinct problems with different objectives and different
solutions. Class imbalance refers to the setting in which training class
frequencies are unequal; the standard remedies, oversampling, undersampling,
and frequency-derived loss reweighting, aim to equalize each class's influence
on the learned decision boundary. Cost-sensitive classification, by contrast,
addresses the setting in which different misclassification errors carry different
real-world penalties. These penalties can be encoded in an $N \times N$ cost matrix $\mathbf{C}$ whose
entry $C[k,j]$ specifies the cost incurred when a sample of true class~$j$ is
predicted as class~$k$~\cite{elkan2001}. These costs are determined by domain
knowledge (clinical severity, ecological consequence, economic impact) and
need not bear any relationship to class frequencies. The two problems can
co-occur, but addressing one does not address the other: a class-balanced
dataset still requires cost-sensitive methods if certain misclassification pairs
are more consequential than others~\cite{elkan2001,ling2008costsensitive}.

The theoretical foundations for cost-sensitive classification are well
established. For binary classification, Elkan~\cite{elkan2001} proved that the
Bayes-optimal decision threshold under an arbitrary $2 \times 2$ cost matrix is
%
\begin{equation}
\label{eq:elkan_threshold}
p^{*} = \frac{C_{\mathrm{FP}} - C_{\mathrm{TN}}}
            {C_{\mathrm{FP}} - C_{\mathrm{TN}} + C_{\mathrm{FN}} - C_{\mathrm{TP}}},
\end{equation}
%
where $C_{\mathrm{FP}}$, $C_{\mathrm{TN}}$, $C_{\mathrm{FN}}$, and
$C_{\mathrm{TP}}$ denote the costs of false positive, true negative, false
negative, and true positive decisions, respectively. This result shows that any
economically coherent binary cost matrix reduces to a single degree of freedom
controlling the decision threshold. For the multiclass setting with $N$ classes,
Bayes decision theory prescribes the cost-minimizing prediction rule
%
\begin{equation}
\label{eq:bayes_cost_rule}
\hat{y}(\mathbf{x}) = \operatorname*{arg\,min}_{k \in \{1,\dots,N\}}
  \sum_{j=1}^{N} C[k,j] \cdot P(Y = j \mid \mathbf{x}),
\end{equation}
%
which selects the class that minimizes expected misclassification cost given the
posterior class probabilities. This rule requires both a well-calibrated
probabilistic classifier and an explicit cost matrix, and is the starting point
for all cost-sensitive
methods~\cite{elkan2001,domingos1999metacost}. Domingos~\cite{domingos1999metacost}
operationalized this principle in MetaCost, a wrapper method that uses bagging
to estimate class probabilities, then relabels each training example with the
class minimizing expected cost under an arbitrary $N \times N$ matrix. MetaCost
handles the full multiclass cost-sensitive setting, yet no direct deep learning
adaptation exists.

Dmochowski et al.~\cite{dmochowski2010jmlr} provided the key theoretical bridge
to loss-function-based approaches by proving that the cost-weighted negative
log-likelihood
%
\begin{equation}
\label{eq:cost_weighted_ce}
\mathcal{L}_{\mathrm{CW}} = -\frac{1}{n}\sum_{i=1}^{n}
  w_{y_i} \log p_{y_i}(\mathbf{x}_i)
\end{equation}
%
is a tight, convex upper bound on the empirical misclassification cost, where
$w_{y_i}$ is a per-class weight derived from the cost matrix. This result
justifies cost-weighted cross-entropy as a surrogate loss for cost-sensitive
training. However, Dmochowski et al.\ also showed that under model
misspecification, the risk-minimizing solution varies with the cost ratio,
meaning that a single set of learned weights may not be optimal across all
operating points.

Despite these mature theoretical foundations, the translation to deep learning
has been limited. Chung, Lin, and Yang~\cite{chung2016csdnn} published the first
work addressing cost-sensitive deep learning with explicit $N \times N$ cost
matrices. Their CSDNN framework introduced the SOSR (Softmax-based Output and
Sensitivity-aware Regularization) loss, which embeds cost information directly
into training, and combined it with cost-aware autoencoders for pre-training.
Evaluated on MNIST, CIFAR-10, USPS, and SVHN with randomized proportional cost
matrices, CSDNN demonstrated that incorporating cost information during training
outperforms approaches that apply cost-sensitive decisions only at prediction
time. The cost matrices used were arbitrary and not derived from class
frequencies, making this work genuinely cost-sensitive rather than a
class-imbalance rebalancing method.

Galdran et al.~\cite{galdran2020} proposed a cost-sensitive regularization
approach applied to diabetic retinopathy grading. Rather than replacing the
standard training loss, they introduced an additive cost-sensitive regularizer
%
\begin{equation}
\label{eq:cs_regularizer}
\mathcal{L}_{\mathrm{CS}}(\mathbf{x}, y) = \sum_{k=1}^{N}
  M[y, k] \cdot p_k(\mathbf{x}),
\end{equation}
%
where $M$ is an $N \times N$ cost matrix and $p_k(\mathbf{x})$ is the predicted
probability for class~$k$. This term penalizes the model for assigning
probability mass to classes whose confusion with the true class~$y$ is costly.
The total training loss combines a standard objective with the regularizer:
%
\begin{equation}
\label{eq:cs_total_loss}
\mathcal{L}_{\mathrm{total}} = \mathcal{L}_{\mathrm{base}}
  + \lambda \cdot \mathcal{L}_{\mathrm{CS}},
\end{equation}
%
where $\lambda$ controls the regularization strength. A key empirical finding
was that using the cost-sensitive term as the sole training loss causes CNNs to
collapse to trivial solutions, predicting a single class for all inputs. This
training instability necessitates the regularization formulation. Applied with
distance-based cost matrices ($M[i,j] = |i - j|$ for L1 and
$M[i,j] = (i-j)^2$ for L2) reflecting disease severity, the approach achieved
a 3--5\% improvement in quadratic-weighted kappa on the EyePACS diabetic
retinopathy dataset~\cite{galdran2020}. This work remains one of very few
open-source PyTorch implementations supporting arbitrary $N \times N$ cost
matrices.

A result with significant implications for this research was established by Chen
et al.~\cite{chen2025csada}, who identified a fundamental obstacle to
cost-sensitive training in modern deep networks. They proved that when
overparameterized neural networks interpolate training data (achieving zero
training loss), cost-weighted empirical risk minimization collapses to standard
ERM, because the loss is zero regardless of the cost weights applied. Formally,
if $\mathcal{L}_{\mathrm{CW}}(\theta^{*}) = 0$ for parameters $\theta^{*}$
that interpolate the data, then the gradient
$\nabla_\theta \mathcal{L}_{\mathrm{CW}}(\theta^{*}) = \mathbf{0}$ irrespective
of the values of the cost weights~$w_y$. This result undermines cost-weighted
cross-entropy, per-class loss reweighting, and any other approach that modifies
the loss by multiplicative cost factors when applied to models with sufficient
capacity to memorize the training set. To address this, Chen et al.\ proposed
CSADA (Cost-Sensitive Adversarial Data Augmentation), a bi-level optimization
framework in which the inner loop generates adversarial examples targeted at
high-cost misclassification pairs via multi-step gradient ascent, and the outer
loop trains on both the original and adversarial examples with cost-aware
penalties. Evaluated on MNIST, CIFAR-10, and a pharmacy medication image dataset
with expert-assigned pairwise costs, CSADA maintained overall accuracy while
reducing average test cost by 27--48\% relative to standard
training~\cite{chen2025csada}.

Several specific gaps define the research agenda for cost-sensitive deep learning
with modern architectures. No Vision Transformer, ConvNeXt, or EfficientNet
architecture has been paired with explicit $N \times N$ cost matrices, and no SSL (self-supervised learning) foundation model backbone, such as DINOv3, CLIP, or their successors, has been
fine-tuned with a cost-sensitive objective. No published work combines
cost-sensitive regularization with logit adjustment, and no work
has attempted to extend logit adjustment from class priors to arbitrary cost
matrices. Finally, the overparameterization result of Chen et al.\ suggests that
parameter-efficient fine-tuning methods, which do not interpolate the training
data, may be essential for cost matrices to meaningfully influence the learned
decision boundary, but this interaction has not been investigated. The practical
consequence is that a domain expert, a clinician who knows that undergrading
proliferative diabetic retinopathy is far more dangerous than overgrading mild
disease, or an ecologist who knows that confusing an invasive species with a
native species has different management consequences than the reverse, currently
has no existing tool to encode that knowledge into a modern deep learning
training pipeline.

\section{Proposed Study}
\label{sec:proposed_study}

\subsection{Overview of Proposed Approach}
\label{sec:overview}

The proposed research consists of two systems that address complementary stages
of the applied image classification pipeline. The first system, FLOW-DC
(Flexible Large-scale Orchestrated Workflow for Data
Collection)~\cite{deng2025flowdc}, is a distributed image downloading framework
that constructs ML-ready datasets from URL manifests hosted across heterogeneous
web servers. The second system, NICME model (Non-Identical Matrix Embedding), is a
cost-sensitive image classification framework that trains deep learning models
using explicit $N \times N$ misclassification cost matrices on modern vision
backbones. The two systems are connected by a shared workflow: FLOW-DC produces
the training datasets on which NICME model operates, and the class-distribution
fidelity of the downloaded dataset directly affects the validity of the
cost-sensitive training that follows.

FLOW-DC extends the manager-worker paradigm of
TaskVine~\cite{slydelgado2023taskvine} with an application-layer congestion
control algorithm called PAARC (Policy-Aware Adaptive Request
Controller)~\cite{deng2025flowdc}. The manager machine partitions a dataset manifest into fixed-length chunks using a host-dispersed partitioning algorithm that separates URLs from the same host across distinct chunks, then assigns these chunks to worker machines as TaskVine tasks.

NICME model is a PyTorch-based image classification framework that extends the
standard training pipeline with support for user-defined cost matrices.
The proposed research will extend this framework in two directions. First, we will integrate
the foundation model backbone
DINOv3~\cite{darcet2025dinov3} with ViT and ConvNeXt using parameter-efficient fine-tuning via
LoRA~\cite{hu2022lora}.
Second, we will develop and evaluate the combined loss function that integrates a cost-sensitive
regularizer~\cite{galdran2020} and logit adjustment~\cite{menon2021logitadjust}, adapting the logit adjustment term to
operate with cost-matrix-derived offsets rather than class priors.

\subsection{FLOW-DC: Distributed Downloading with PAARC}
\label{sec:flowdc}

FLOW-DC addresses Objective~1 through three interacting components: a
host-aware dataset partitioning algorithm, the PAARC congestion control
algorithm operating on each worker, and the TaskVine-based distributed
orchestration layer. This section details each component. The full system
design and implementation are described in Deng et
al.~\cite{deng2025flowdc}.

\paragraph{Host-Aware Partitioning.}
Before distributing work to workers, the manager partitions the dataset
manifest into chunks that disperse URLs from the same host server across
distinct partitions. This dispersion is important because it prevents any
single worker from concentrating requests on one host, distributing the
load across the worker pool and reducing the likelihood of server-side
throttling. Let
$H = \{h_1, h_2, \ldots, h_m\}$ denote the set of $m$ unique hosts extracted
from the manifest, where each host $h_i$ has URL count $c_i$ and URL set
$U_i = \{u_{i,1}, u_{i,2}, \ldots, u_{i,c_i}\}$. Let $k$ denote the desired
number of partitions (a hyperparameter). For each host, the algorithm assigns its URLs to
partitions via offset round-robin allocation, where the starting partition
rotates per host to balance the distribution of remainder URLs:
%
\begin{equation}
\label{eq:host_dispersion}
\forall\; u_{i,j} \in U_i:\quad
u_{i,j} \longrightarrow P_{((j - 1 + \delta_i) \bmod k) + 1},
\quad \delta_i = (i-1) \bmod k.
\end{equation}
%
The per-host offset $\delta_i$ rotates the round-robin starting partition for
each successive host, ensuring that the remainder URLs (those arising when
$c_i$ is not evenly divisible by $k$) are spread across different partitions
rather than accumulating in the same low-indexed partitions. This provides two
simultaneous guarantees: each partition receives exactly
$\lfloor c_i / k \rfloor$ or $\lceil c_i / k \rceil$ URLs from any single
host $h_i$, providing optimal per-host dispersion, and global partition sizes
remain balanced to within at most $\lfloor m / k \rfloor$ URLs of each other
in the worst case. When $c_i \geq k$, every partition receives at least one
URL from host $h_i$, ensuring no single worker bears a disproportionate share
of requests to any server.

\paragraph{PAARC Algorithm.}
The Policy-Aware Adaptive Request Controller operates on each worker node and
dynamically adjusts the number of concurrent HTTP requests issued to each host
server. The primary control variable is the concurrency $C$, the number of
simultaneous in-flight requests. The download rate
$R$ emerges naturally from Little's Law, $R = C / T$, where $T$ is the mean
response time~\cite{deng2025flowdc}. This concurrency-primary design is
self-regulating: when server latency increases, throughput decreases
automatically without requiring an explicit rate adjustment.

PAARC monitors time-to-first-byte (TTFB) as its primary latency signal. For
each control interval, the controller computes the 10th, 50th, and 95th
percentiles of successful request TTFBs and smooths them with an exponential
moving average:
%
\begin{equation}
\label{eq:ema}
p_k^{(t)} = \alpha \cdot p_k^{\mathrm{raw}}
           + (1 - \alpha) \cdot p_k^{(t-1)},
\end{equation}
%
where $p_k$ denotes the $k$-th percentile and $\alpha=0.3$ is a hyperparameter controlling the smoothing
factor. The baseline latency RTprop is defined as the minimum observed
10th-percentile TTFB over a sliding window:
%
\begin{equation}
\label{eq:rtprop}
\mathrm{RTprop} = \min_{t' \in [t - W,\, t]} p_{10}^{(t')},
\end{equation}
%
where $W = 10$~seconds is a hyperparameter. Latency degradation is detected when median or tail
latencies exceed threshold multiples of RTprop:
%
\begin{equation}
\label{eq:degradation}
\mathrm{degraded} = \bigl(p_{50} > \mathrm{RTprop} \times \theta_{50}\bigr)
  \;\lor\;
  \bigl(p_{95} > \mathrm{RTprop} \times \theta_{95}\bigr),
\end{equation}
%
with hyperparameters $\theta_{50} = 1.5$ and $\theta_{95} = 2.0$.

\paragraph{PAARC State Machine.}
PAARC implements a five-state finite state machine governing concurrency
adjustment decisions~\cite{deng2025flowdc}. In the \textbf{INIT} state, the
controller operates at the hyperparameter $C_{\mathrm{init}}$ concurrent requests and collects
initial TTFB samples to establish the RTprop baseline. Upon collecting
sufficient samples, the controller transitions to \textbf{STARTUP}, where it
discovers the maximum sustainable concurrency through additive growth. Each
control interval, the concurrency is increased by a fixed hyperparameter increment $\Delta C_{\mathrm{startup}}$. STARTUP
terminates when either an overload signal occurs (transitioning to BACKOFF) or a
latency-based plateau is detected; that is, when the latency degradation
condition in Equation~\ref{eq:degradation} is triggered using elevated
hyperparameter thresholds $\theta_{50}^{\mathrm{startup}}$ and
$\theta_{95}^{\mathrm{startup}}$. Upon plateau detection, PAARC sets the
concurrency ceiling $C_{\mathrm{ceiling}}$ to the current concurrency and
computes the operating point:
%
\begin{equation}
\label{eq:operating}
C_{\mathrm{operating}} = \lfloor C_{\mathrm{ceiling}} \times \mu \rfloor,
\end{equation}
%
where $\mu = 0.75$ is the utilization factor (a hyperparameter), then transitions to the
steady-state \textbf{PROBE\_BW} phase. In PROBE\_BW, the controller maintains
concurrency at or below $C_{\mathrm{operating}}$ and periodically probes for
additional capacity via conservative additive increase. If an overload signal
occurs, the ceiling is revised downward by a multiplicative decrease hyperparameter
$\beta = 0.5$ and the controller transitions to \textbf{BACKOFF}. The
controller periodically enters \textbf{PROBE\_RTT} to refresh the RTprop
estimate by reducing concurrency by a fraction $\rho$ of the current value:
%
\begin{equation}
\label{eq:probe_rtt}
C_{\mathrm{probe}} = \max\bigl(C_{\mathrm{min}},\;
  \lfloor C \times \rho \rfloor\bigr),
\end{equation}
%
where $\rho = 0.5$ and $C_{\mathrm{min}} = 1$ are hyperparameters,
collecting fresh latency samples at reduced load, then gradually restoring
concurrency over multiple steps to avoid burst-induced congestion. BACKOFF
handles recovery from overload conditions (HTTP~429, 5xx, or connection-level
errors) by enforcing a cooldown period before restoring concurrency to
$C_{\mathrm{operating}}$ and returning to PROBE\_BW. This design draws on
BBR's model-based approach to congestion
control~\cite{cardwell2017bbr} but adapts its mechanisms for the application
layer, where congestion signals are HTTP status codes and TTFB measurements
rather than packet-level acknowledgements.

\paragraph{Distributed Execution.}
The manager instantiates a TaskVine workflow in which each partition is
submitted as an independent task with declared input files (partition parquet,
download scripts, configuration JSON) and output files (\eg compressed tar archive,
JSON summary). Workers connect to the manager utilizing the manager's IP and receive task
assignments from the pending queue. TaskVine handles input file transfer, task
scheduling, and result collection. Workers that disconnect mid-workflow do not
cause data loss; TaskVine detects the disconnection and re-queues incomplete
tasks for available workers. This elastic participation model requires only
network reachability between the manager and workers: no shared filesystem, no
common job scheduler, and no synchronized start
times~\cite{slydelgado2023taskvine}.

\subsection{NICME: Cost-Sensitive Learning using Logit Adjustment and Regularization}
\label{sec:loss_function}

The core algorithmic contribution of NICME model is a loss function that
integrates two cost-sensitive mechanisms, logit adjustment and cost-sensitive
regularization, into a single training objective. Each mechanism addresses a
different aspect of cost-aware learning.

The first component is a cost-matrix-driven logit adjustment. For each training
sample, the model's predicted class
$\hat{k} = \arg\max_k f_k(\mathbf{x})$ is compared to the true class~$y$. For
misclassified samples ($\hat{k} \neq y$), the logit corresponding to the
predicted class is increased proportionally to both the cost of that specific
misclassification and the logit gap between the predicted and true classes:
%
\begin{equation}
\label{eq:logit_adjust_cost}
\tilde{f}_{\hat{k}}(\mathbf{x}) = f_{\hat{k}}(\mathbf{x})
  + C[y, \hat{k}] \cdot \bigl| f_{\hat{k}}(\mathbf{x})
  - f_y(\mathbf{x}) \bigr|,
\end{equation}
%
where $C[y, \hat{k}]$ is the entry of the user-defined cost matrix
corresponding to misclassifying true class~$y$ as class~$\hat{k}$, and
$f_k(\mathbf{x})$ denotes the logit for class~$k$. For correctly classified
samples, the logits are unchanged. This adjustment amplifies the cross-entropy
gradient for high-cost misclassifications: when $C[y, \hat{k}]$ is large, the
modified logit $\tilde{f}_{\hat{k}}$ makes the model appear more confident in
the costly wrong prediction, producing a stronger corrective gradient during
backpropagation. The cross-entropy loss is then computed on the adjusted logits:
%
\begin{equation}
\label{eq:ce_adjusted}
\mathcal{L}_{\mathrm{LA}} = -\frac{1}{n} \sum_{i=1}^{n}
  \log \frac{\exp(\tilde{f}_{y_i}(\mathbf{x}_i))}
            {\sum_{k=1}^{N} \exp(\tilde{f}_k(\mathbf{x}_i))}.
\end{equation}

The second component is a cost-sensitive regularizer introduced by Galdran et
al.~\cite{galdran2020}, which penalizes the model for assigning probability
mass to classes whose confusion with the true class is costly. This term
operates on the original (non-adjusted) logits to ensure the two gradient paths
are independent:
%
\begin{equation}
\label{eq:cs_reg_nicme}
\mathcal{L}_{\mathrm{CS}} = \frac{1}{n} \sum_{i=1}^{n} \sum_{k=1}^{N}
  \bar{M}[y_i, k] \cdot \sigma_k(\mathbf{x}_i),
\end{equation}
%
where $\bar{M}$ is the cost matrix normalized to $[0, 1]$ and
$\sigma_k(\mathbf{x}_i) = \mathrm{softmax}(f(\mathbf{x}_i))_k$ is the
predicted probability for class~$k$ computed from the original logits. The
normalization prevents unbounded penalty magnitudes when cost matrix entries
span several orders of magnitude.

The combined loss function is:
%
\begin{equation}
\label{eq:combined_loss}
\mathcal{L}_{\mathrm{NICME}} = \mathcal{L}_{\mathrm{LA}}
  + \lambda_{\mathrm{eff}} \cdot \mathcal{L}_{\mathrm{CS}},
\end{equation}
%
where $\lambda_{\mathrm{eff}}$ is the effective regularization weight (a hyperparameter). To
prevent the regularization term from destabilizing early training, a failure
mode identified by Galdran et al.~\cite{galdran2020}, who found that
cost-sensitive losses alone cause CNNs to collapse to trivial solutions, we
apply a linear warmup schedule:
%
\begin{equation}
\label{eq:warmup}
\lambda_{\mathrm{eff}}(e) = \lambda \cdot \min\!\left(1,\;
  \frac{e}{E_{\mathrm{warmup}}}\right),
\end{equation}
%
where $e$ is the current epoch and the hyperparameters $E_{\mathrm{warmup}}$
and $\lambda$ specify the warmup duration in epochs and the target
regularization strength, respectively. During warmup, the
model trains primarily on $\mathcal{L}_{\mathrm{LA}}$; the regularizer's
contribution increases linearly until reaching full strength at epoch
$E_{\mathrm{warmup}}$.

\paragraph{Backbone Architecture and Parameter-Efficient Fine-Tuning.} 
NICME model currently implements ResNet-50~\cite{he2016resnet} and
ConvNeXt~\cite{liu2022convnext} backbones using custom base classes that
decouple the model architecture from HuggingFace's transformers library. Both
architectures load ImageNet-pretrained weights and replace the final
classification head with a linear layer matching the number of target classes.
Stage freezing is supported: setting the hyperparameter~$s$ (number of frozen stages) freezes
all parameters in the first $s$ convolutional stages, restricting gradient
updates to the later stages and the classification head.

The proposed extension integrates DINOv3~\cite{darcet2025dinov3} pretraining combined with ConvNeXt and ViT as backbone architectures via
parameter-efficient fine-tuning using LoRA (Low-Rank
Adaptation)~\cite{hu2022lora}. LoRA constrains the weight update to a low-rank
decomposition: for a pretrained weight matrix
$\mathbf{W}_0 \in \mathbb{R}^{d \times d}$, the adapted weight is
%
\begin{equation}
\label{eq:lora}
\mathbf{W} = \mathbf{W}_0 + \mathbf{B}\mathbf{A},
\end{equation}
%
where $\mathbf{B} \in \mathbb{R}^{d \times r}$ and
$\mathbf{A} \in \mathbb{R}^{r \times d}$ with rank $r \ll d$. Only
$\mathbf{A}$ and $\mathbf{B}$ are trainable; $\mathbf{W}_0$ remains frozen.
This reduces the number of trainable parameters to 0.1--2\% of the full model,
which provides two advantages for cost-sensitive learning. First, it avoids the
overparameterization collapse identified by Chen et
al.~\cite{chen2025csada}: because the vast majority of parameters are frozen,
the model cannot fully interpolate the training data, and the cost weights in
$\mathcal{L}_{\mathrm{NICME}}$ retain their influence on the learned decision
boundary. Second, the LIFT result~\cite{tian2024lift} demonstrated that
lightweight fine-tuning with logit-adjusted loss preserves the class-conditional
feature distributions learned during pretraining, which heavy fine-tuning
destroys. LoRA applied to the linear projection layers of DINOv3 and
ConvNeXt, with rank $r \in \{4, 8, 16\}$ as a hyperparameter, is the
proposed default configuration.

Post-hoc temperature scaling will be applied as a final calibration step. Given
a trained model with logits $f(\mathbf{x})$, the calibrated probabilities are
$\sigma(f(\mathbf{x}) / T)$ where the temperature $T > 0$ is optimized on a
held-out validation set by minimizing the negative
log-likelihood~\cite{guo2017calibration}.

The primary metric used to evaluate whether the trained model satisfies the
user-specified cost preferences is the average test cost (ATC), defined as the
mean per-sample cost incurred by the model's predictions on the test set:
\begin{equation}
\label{eq:atc}
\mathrm{ATC} = \frac{1}{n} \sum_{i=1}^{n} C[\hat{y}_i, y_i],
\end{equation}
%
where $n$ is the number of test samples, $\hat{y}_i$ is the model's prediction
for sample $i$, $y_i$ is the true class, and $C[\hat{y}_i, y_i]$ is the
corresponding entry of the user-specified cost matrix. ATC is the quantity
that hyperparameter optimization and final model selection will minimize.
Alongside ATC, we will additionally track per-class recall stratified by cost
category (high, medium, low), and overall classification accuracy, to confirm that cost reductions are not achieved by
degrading general classification performance or calibration. 

\subsection{Prior Work and Preliminary Implementations}
\label{sec:preliminary}

The research proposed in this document is inspired by our work AraNet, a two-stage convolutional neural network for spider species
classification~\cite{deng2026aranet}. The lessons from AraNet directly motivated both FLOW-DC and NICME model.

\paragraph{AraNet: Spider Species Classification via Multi-Phase Deep Learning.}
AraNet was developed to address the problem of automated spider species
identification, a task with applications in habitat monitoring and emergency
medical treatment~\cite{deng2026aranet}. Two practical challenges encountered during the AraNet project directly
motivated the present research.

The first was dataset construction. Assembling
the iNatSpiders1k training set required downloading 548,689 images from
iNaturalist's AWS-hosted repository. While iNaturalist serves images from a
single high-capacity S3 endpoint, expanding the dataset to utilize other biodiversity
data sources, particularly GBIF, which aggregates records across over 500
institutional servers with varying capacity and rate-limiting
policies, revealed the limitations of existing download tools. This experience
motivated the development of FLOW-DC and its PAARC congestion control algorithm.

The second challenge was the recognition that standard accuracy-maximizing
training does not capture the asymmetric consequences of misclassification in
safety-relevant domains. In the Australian spider classification task,
misidentifying a venomous species such as the Sydney funnel-web
(\textit{Atrax robustus}) as a non-dangerous species could delay life-saving
medical intervention, while the reverse error (classifying a harmless species
as dangerous) results in unnecessary caution but no physical harm.
This observation motivated the development of NICME's cost-matrix-based training framework and directly
informs the design of the combined loss function described in
Section~\ref{sec:loss_function}.

\paragraph{FLOW-DC: Production Deployment for Biotrove Dataset}
FLOW-DC has been deployed at production scale for the construction of the
Biotrove dataset, a large-scale biodiversity image collection contributed to the
NeurIPS 2024 Datasets and Benchmarks Track~\cite{yang2024biotrove}. In this
deployment, FLOW-DC downloaded 38.7~million images totaling 13~TiB from GBIF's
network of institutional servers using 10 Jetstream2 cloud workers. The core architecture of FLOW-DC including the PAARC algorithm has been constructed~\cite{deng2025flowdc}.

\paragraph{NICME Preliminary Model}
The prototype NICME model framework has been implemented in PyTorch using the HuggingFace
Trainer API. The framework currently supports ResNet-50~\cite{he2016resnet} and
ConvNeXt~\cite{liu2022convnext} backbones with configurable stage freezing, and
utilizes an augmented cross-entropy loss that integrates logit adjustment with cost matrix values. The
cost matrix is accepted as a user-defined input and propagated throughout model training.

\section{Proposed Experiments}
\label{sec:experiments}

The experimental plan is organized around the two research objectives stated in
Section~\ref{sec:objectives}. For Objective~1 (FLOW-DC), the evaluation
focuses on rate-limit compliance, server-fault resilience, and scaling behavior rather than raw download
throughput, since both FLOW-DC and its primary baseline are ultimately
bottlenecked by available network bandwidth. For Objective~2 (NICME model), the
evaluation focuses on whether the proposed cost-sensitive loss function lowers
the average test cost (ATC) induced by a user-specified cost matrix, while
maintaining overall classification accuracy. Each objective is evaluated against a set of
baselines selected to isolate the contribution of specific design decisions.

\subsection{Baseline Comparisons}
\label{sec:baselines}

\paragraph{FLOW-DC Baselines.}
The primary baseline to evaluate FLOW-DC is \textbf{img2dataset}~\cite{beaumont_img2dataset}, the
most widely deployed open-source tool for large-scale image downloading. As
detailed in Section~\ref{sec:flowdc}, img2dataset uses synchronous thread-based
HTTP via \texttt{urllib.request} with no per-host congestion control, treats
HTTP~429 responses as permanent failures, and requires Apache PySpark for
distributed execution. We believe img2dataset is the appropriate primary baseline because
it represents the current standard against which any new downloading tool must
be measured, and its published benchmarks (approximately 90 samples per second
per CPU core, scaling to roughly 10,000 images per second across a 10-node
Spark cluster~\cite{schuhmann2022laion5b}) provide the throughput reference
for establishing parity.

\paragraph{NICME Model Baselines.}
Three baselines are used to evaluate NICME model's cost-sensitive loss function,
each isolating a different aspect of the proposed approach. All baselines will
be evaluated using the same backbone architecture, dataset, and training
hyperparameters to ensure that differences in performance are attributable to
the loss function rather than to confounding factors.

The first baseline is \textbf{standard cross-entropy} (CE), which assigns
equal penalty to all misclassification errors:
%
\begin{equation}
\label{eq:baseline_ce}
\mathcal{L}_{\mathrm{CE}} = -\frac{1}{n} \sum_{i=1}^{n}
  \log \frac{\exp(f_{y_i}(\mathbf{x}_i))}
            {\sum_{k=1}^{N} \exp(f_k(\mathbf{x}_i))}.
\end{equation}
%
This is the cost-insensitive baseline against which any cost-sensitive method
must demonstrate improvement in average test cost (ATC) as defined in
Equation~\ref{eq:atc}. Standard CE is expected to minimize overall error rate
but not to preferentially avoid high-cost misclassification pairs.

The second baseline is \textbf{cost-sensitive regularization alone} (NULS-CS),
as proposed by Galdran et al.~\cite{galdran2020}. This uses the loss
$\mathcal{L}_{\mathrm{base}} + \lambda \cdot \mathcal{L}_{\mathrm{CS}}$ with
the regularizer defined in Equation~\ref{eq:cs_regularizer} but without the
logit adjustment component. NULS-CS achieved a 3--5\% improvement in
quadratic-weighted kappa on diabetic retinopathy grading using a ResNeXt-50
backbone~\cite{galdran2020}. This baseline isolates the contribution of the
logit adjustment term: if NICME model outperforms NULS-CS, the improvement is
attributable to the cost-matrix-driven logit adjustment described in
Section~\ref{sec:loss_function}.

The third baseline is \textbf{logit-adjusted cross-entropy} as proposed by
Menon et al.~\cite{menon2021logitadjust}, which adds the log class prior
$\log \pi_y$ to each logit. Although
this method was designed for long-tailed class distributions rather than
cost-sensitive classification, it is the closest published method to
NICME model's logit adjustment mechanism and serves as a reference for the
contribution of replacing class-prior offsets with cost-matrix-derived
adjustments. On balanced datasets where $\pi_y = 1/N$ for all classes, this
baseline reduces to standard CE, providing an additional control.

\subsection{Experimental Methodology}
\label{sec:methodology}

\paragraph{FLOW-DC Experiments.}
FLOW-DC will be evaluated using two differing manifests that
separate content delivery network (CDN) based workloads from server-diverse workloads. The first
manifest is a CDN-dominated control drawn from DataComp
small~\cite{gadre2023datacomp}, which contains 12.8~million image-text pairs
derived from Common Crawl. Because Common Crawl-sourced URLs terminate
predominantly on high-capacity CDN infrastructure, this manifest tests
throughput parity between FLOW-DC and img2dataset under conditions where
congestion control provides low benefit; neither tool should trigger rate
limiting. The second manifest is a server-diverse experimental condition
constructed by filtering the GBIF occurrence snapshot to a subset of
approximately 1--3~million URLs spanning 50--100 distinct institutional
servers. This manifest tests PAARC's core value
proposition: adaptive throttling on capacity-limited servers.

Both FLOW-DC (with PAARC enabled), and
img2dataset will be run on both manifests.
For single-machine experiments, the primary platform is the University of
Arizona HPC Puma cluster. For distributed
experiments, 10 Jetstream2 cloud instances will be used (matching the Biotrove dataset generation
configuration~\cite{deng2025flowdc}). Scaling efficiency will be measured at
$N = 1, 2, 5, 10$ workers.

\paragraph{FLOW-DC Metrics.}
The primary throughput metric is the aggregate download rate
$T_{\mathrm{agg}} = N_{\mathrm{success}} / t_{\mathrm{download}}$ (images per
second), where $t_{\mathrm{download}}$ excludes partitioning, packaging, and
transfer time. Both images-per-second and TiB-per-hour will be reported, with
TiB-per-hour preferred for cross-study comparison because image sizes differ
across datasets. The primary metric for evaluating rate-limit compliance and
server-fault resilience is the per-host server-side failure rate (SFR),
which aggregates all responses indicative of server throttling or overload:
HTTP~429 (Too Many Requests), 5xx server errors (500, 502, 503, 504),
connection resets, and request timeouts:
%
\begin{equation}
\label{eq:sfr}
\mathrm{SFR}(h) = \frac{N_{429}(h) + N_{5xx}(h) + N_{\mathrm{timeout}}(h)
  + N_{\mathrm{reset}}(h)}{N_{\mathrm{total}}(h)} \times 100\%.
\end{equation}
%
Scaling efficiency is
defined as:
%
\begin{equation}
\label{eq:scaling}
E(N) = \frac{T_{\mathrm{agg}}(N)}{N \cdot T_{\mathrm{agg}}(1)},
\end{equation}
%
where $T_{\mathrm{agg}}(N)$ is the aggregate throughput with $N$ workers.
Output dataset quality is measured by the class distribution fidelity (CDF):
%
\begin{equation}
\label{eq:cdf}
\mathrm{CDF} = 1 - D_{\mathrm{KL}}\!\left(P_{\mathrm{downloaded}}
  \;\|\; P_{\mathrm{requested}}\right),
\end{equation}
%
where $P_{\mathrm{downloaded}}$ and $P_{\mathrm{requested}}$ are the class
frequency distributions of successfully downloaded images versus the original
manifest.

\paragraph{NICME Model Experiments.}
NICME model will be evaluated on four datasets of increasing complexity,
each selected to provide a domain-motivated cost matrix. The first is a binary spider classification
dataset (2 classes, 500 images per class) distinguishing black widows
(\textit{Latrodectus}) from false widows (\textit{Steatoda}), two genera that
are frequently confused due to similar morphology but carry vastly different
medical consequences. The $2 \times 2$ cost matrix assigns high cost to
misclassifying a black widow as a false widow (risking delayed antivenom
treatment) and low cost to the reverse conservative error. This dataset serves
as the simplest baseline for validating the cost-sensitive loss function in a
controlled binary setting. The second is the AusSpiders dataset (9 classes,
64,291 images)~\cite{deng2026aranet}, which extends the binary venomous
versus non-venomous cost structure to a multiclass setting. The class
distribution will be balanced by downsampling to equalize class frequencies,
isolating the effect of the cost matrix from class-imbalance confounds. The
third is the EyePACS diabetic retinopathy grading
dataset~\cite{galdran2020}, which provides a 5-class ordinal structure (grades
0--4) with clinically motivated asymmetric costs: undergrading severe disease
(e.g., predicting grade~0 for a grade~4 sample) carries higher cost than
overgrading mild disease. This dataset enables direct comparison with the
NULS-CS baseline of Galdran et al.~\cite{galdran2020}, who reported QWK scores
on the same data. The fourth is the Pharmacy Medication Image (PMI) dataset
(20 classes, 13,955 images)~\cite{chen2025csada}, which contains top-down pill
images identified by National Drug Code (NDC). Expert costs were assigned by a
pharmacist to four critical medication-confusion pairs based on clinical
danger; for example, confusing Amiodarone (a cardiac drug) with Allopurinol (a gout
medication) risks pulmonary toxicity. Non-critical pairs are assigned
unit cost. This dataset provides the most complex cost structure and enables
direct comparison with the CSADA baseline of Chen et
al.~\cite{chen2025csada}.

For each dataset, all baselines and the proposed
$\mathcal{L}_{\mathrm{NICME}}$ loss will be trained using the same backbone, the same data splits (80/10/10 train/validation/test, stratified by
class), and the same training hyperparameters (AdamW optimizer, cosine
annealing, early stopping on validation loss). To ensure statistical
reliability, each configuration will be run across a minimum of 5 random seeds.

Optimal hyperparameters will be determined via Bayesian hyperparameter
optimization (HPO) using Optuna~\cite{akiba2019optuna} with the Tree-structured
Parzen Estimator (TPE) sampler. For each dataset and backbone combination,
Optuna will search over the cost-sensitive hyperparameters introduced in
Section~\ref{sec:loss_function} (the regularization strength~$\lambda$, the
warmup duration~$E_{\mathrm{warmup}}$, and the LoRA rank~$r$) as well as
standard training hyperparameters including learning rate, weight decay, and
batch size. The optimization objective is to minimize ATC on the validation
split, ensuring that the selected hyperparameters directly optimize the
cost-sensitive goal rather than overall accuracy.

\paragraph{NICME Model Metrics.}
The primary metric is the average test cost (ATC) as defined in
Equation~\ref{eq:atc}. The cost reduction ratio (CRR) will be reported as:
%
\begin{equation}
\label{eq:crr}
\mathrm{CRR} = \frac{\mathrm{ATC}_{\mathrm{baseline}}
  - \mathrm{ATC}_{\mathrm{proposed}}}
  {\mathrm{ATC}_{\mathrm{baseline}}}.
\end{equation}
%
Per-class recall stratified by cost category (high, medium, low) will be
reported as a secondary metric. Overall classification accuracy and F1-score will also be
tracked to ensure that cost-sensitive training does not trivially improve
high-cost recall by degrading performance on other classes. All results will be
reported as mean $\pm$ standard deviation across seeds, with 95\% confidence
intervals.

\paragraph{Hardware.}
FLOW-DC single-machine experiments will use the UofA HPC Puma cluster (AMD EPYC
7642, 25~Gb/s Ethernet); distributed experiments will use up to 10 Jetstream2
cloud instances (up to 128 cores, 100~Gbps per node). NICME model training will
be conducted on the University of Arizona HPC GPU resources, primarily Nvidia
V100S GPUs (32~GB) and Nvidia A100 GPUs. Local experiments will utilize an Nvidia RTX~5090
GPU (32~GB) for rapid prototyping.

\section{Risks and Alternatives}
\label{sec:risks}

\paragraph{Overparameterization collapse under LoRA.}
The use of LoRA to preserve cost-weight influence has not been empirically
verified for cost-sensitive objectives. If LoRA-adapted models still achieve
sufficiently low training loss that cost weights become ineffective, we will
fall back to investigating cost-sensitive post-hoc decision rules using the Bayes-optimal
rule (Equation~\ref{eq:bayes_cost_rule}), a two-stage approach that requires no modification to the
training loss~\cite{elkan2001}.

\paragraph{Combined loss does not outperform individual components.}
The logit adjustment and regularization terms in
$\mathcal{L}_{\mathrm{NICME}}$ may interfere destructively rather than
complement each other. If no configuration of $\lambda$ and
$E_{\mathrm{warmup}}$ yields improvement over the individual components, the
contribution shifts to optimization and augmentation on the single best performing method between regularization and logit adjustment.

\section{Resources}
\label{sec:resources}

All computational, storage, and software resources for the project are accessible through existing
institutional allocations and personal hardware. Compute time on the
University of Arizona HPC center (Puma CPU nodes and the shared GPU
allocation providing Nvidia V100S and A100 GPUs) is available under the
advisor's standing allocation, and NSF Jetstream2 cloud instances for
distributed FLOW-DC experiments are covered under an existing usage-based
allocation. 

\paragraph{Software.}
All software used in this research is open-source and requires no commercial
licenses. FLOW-DC is implemented in Python using aiohttp (async HTTP),
Polars (data processing), and TaskVine (distributed workflow
management)~\cite{slydelgado2023taskvine}. NICME model is implemented in PyTorch
with the HuggingFace Trainer API, Optuna for hyperparameter
optimization~\cite{akiba2019optuna}, and the timm library for pretrained vision
model access. Version control is managed through GitHub.

\paragraph{Data and Datasets.}
FLOW-DC experiments will use publicly available URL manifests from
GBIF (Parquet format, hosted on AWS Open Data) and DataComp
small~\cite{gadre2023datacomp}. NICME model experiments will use a binary black widow
versus false widow spider dataset (1,000 images), the AusSpiders dataset
(64,291 images, curated from iNaturalist)~\cite{deng2026aranet}, the EyePACS
diabetic retinopathy dataset~\cite{galdran2020}, and the Pharmacy Medication
Image (PMI) dataset (13,955 images)~\cite{chen2025csada}. All datasets will be
publicly accessible.

\section{Proposed Schedule}
\label{sec:schedule}

\begin{table}[H]
    \centering
    \footnotesize
    \renewcommand{\arraystretch}{0.95}
    \caption{Proposed research timeline.}
    \label{tab:schedule}
    \begin{tabular}{@{}lll@{}}
        \toprule
        \textbf{Month} & \textbf{Task} & \textbf{Deliverable} \\
        \midrule
        Spring 2026 & Oral Qualification Exam & Qualification Exam Pass \\
        Spring 2026 & FLOW-DC final throughput experiments & FLOW-DC results \\
        Spring 2026 & FLOW-DC paper draft & FLOW-DC paper \\
        Spring 2026 & NICME binary model experiments & NICME preliminary results \#1 \\
        Summer 2026 & Binary Class NICME Paper & NICME paper \#1 \\
        Summer 2026 & NICME multi-class extension experiments & NICME multiclass results \\
        Fall 2026 & NICME class imbalance methods integration testing & NICME integration results \\
        Fall 2026 & NICME extension paper draft & NICME paper \#2 \\
        Fall 2026 & Dissertation writing & Dissertation paper \\
        Fall 2026 & Defense preparation & Dissertation defense \\
        \bottomrule
    \end{tabular}
\end{table}

\newpage
\begingroup
\small
\setlength{\bibsep}{2pt plus 1pt minus 1pt}
\bibliographystyle{ieeetr}
\bibliography{references}
\endgroup

\end{document}
